\itemhead{第2章 小结与复习 · 第4题（教材p.60）}{XJ-B1-C02-SU-P060-Y1255}
\label{XJ-B1-C02-SU-P060-Y1255}
{\footnotesize 湘教版教材p.60 / PDF第67页 / 第2章 小结与复习 / 第4题}
\field{原题}{从函数观点看一元二次方程和一元二次不等式，可以感悟到数学知识之间的紧密关联：相等与不等是可以相互转化的，正是这种相互转化，使我们可以借助函数图象和方程的解来求解不等式。试结合问题，体会解一元二次不等式的方法与步骤。}
\field{参考答案}{化为标准形式并确定二次项符号；求实根或判别式；画开口与零点草图；按严格或非严格符号选区间；用集合表示，特别检查重根和端点。}
{\small \field{关联知识}{I04 一元二次不等式与参数分类（教材帮p.64、65、66 / PDF73、74、75）；Q01 二次函数、零点、判别式与根系关系（教材帮p.65 / PDF74）}
\field{方法与易错点}{化标准式、查系数、求根或判别式、按图定号，最后与实际范围取交集。 易错：套“内小外大”却忘开口；丢重根或降次；漏整数产量限制。}
\field{答案依据}{答案册未提供本栏目的独立对应答案；依据原题、教材概念补写参考回应。开放性总结不设唯一数值答案。}}
